% !TEX root = ../main.tex
\section{엔도스랭크와 AWP를 나란히 놓고 보기}
\label{ch:appendix-er-awp}

\ref{ch:results}장은 이 논문의 모든 점수를 함께 보고해요. 그리고 각 페이지랭크(PageRank)를, 그 페이지랭크가 매끄럽게 다듬는(평활하는) 원시 차수(화살표를 그냥 센 수)와 견주어 읽어요. 이 부록은 한 층만 걷는 두 점수인 엔도스랭크(EndorseRank)와 AWP(적응형 가중 페이지랭크)를 바로 견준 결과를 모아요. 같은 기간에서의 지표 가족(비슷한 지표 묶음) 평균과 짝지은 차이, 대리 지표(대신 재는 값) 하나하나에 대한 계수, 기간 밖에서의 짝지은 차이, 그리고 모든 트레이더(거래하는 사람) 라벨(나중에 채점에 쓰는 정답 값)에서의 대비예요. 여기 있는 모든 수는 수정 분석(등록 뒤에 금액을 달러로 바꿔 다시 한 분석)에서 나와요. 금액은 다섯 토큰에 걸쳐 달러(USD)로 매겼어요(\ref{sec:amount-revision}절). 등록 분석(미리 등록한 그대로 한 분석)은 부록~\ref{sec:map-registered}에 있어요. 매칭 코호트(모든 점수가 똑같이 점수를 매기는 묶음)나 스펜더(spender, 허락을 받아 대신 쓰는 쪽) 코호트에서 한 비교는 어느 것도 미리 정한 것이 아니에요. 그리고 어느 것도 이 논문의 주장(C)을 담고 있지 않아요. 트레이더 쌍은 등록 분석을 위해 W1과 함께 등록했어요. 여기서는 그 쌍을 고친 점수로 다시 계산해요(\ref{sec:er-awp-neutral-grid}절). 두 점수는 서로 다른 화살표(엣지) 위를 걷고, 서로 다른 구성개념(점수가 재려고 하는 생각)을 위해 만들었어요. 그래서 코호트 전체에서 하는 점검마다, 앞서는 점수는 그 점검이 세는 화살표를 걷는 점수예요. 두 화살표 종류 어느 쪽으로도 만들지 않은 결과는 예외예요. 어느 점수도 그 결과를 위해 만들지 않았기 때문이에요. 그 결과는 \ref{sub:neutral-label-results}절에서 보고해요.

두 점수가 점수를 매기는 범위는 서로 달라요. 매칭 코호트 가운데 스마트 계약(블록체인 위에서 저절로 실행되는 프로그램) $523$개는 허락(allowance, 토큰을 얼마까지 대신 써도 된다는 허락) 그래프(점과 화살표로 이은 그림) 밖에 있고, $2{,}670$개는 송금(transfer) 그래프 밖에 있어요. 또 W0의 스펜더 $17{,}050$개 가운데 $2{,}891$개는 동결(점수를 고정하는 일) 전에 송금 화살표가 하나도 없었어요. 그래서 AWP는 이 계약들에 0점을 줘요. 이와 달리 엔도스랭크는 코호트 계약 쌍을 거의 하나도 동점으로 남기지 않아요(\ref{sec:rank-divergence}절). 기간 밖에서는 두 점수가 자기가 다듬는 원시 차수와 견주어 서로 다른 자리에 서요. AWP는 두 기간 모두에서 나중의 새 송금자(새로 송금을 보내온 주소)를 송금 입차수(송금을 보내온 서로 다른 주소 수)보다 더 잘 줄 세워요. 반면에 엔도스랭크는 두 기간 모두에서 나중의 새 허락 쌍(새로 생긴 주인–스펜더 허락)을 받은 허락 수(허락 입차수)보다 못하게 줄 세워요(\ref{sec:holdout-results}절). 그러므로 엔도스랭크가 AWP보다 앞선다고 해서, 허락을 따라 퍼뜨리는 것(전파)이 허락한 주인(허락을 해 주는 쪽)의 수에 무언가를 보탠다는 것을 보여 주지는 않아요.

\dissertationSubheading[sec:er-awp-same-window]{같은 기간}

표~\ref{tab:alignment-hybrid}은 두 점수의 지표 가족 평균을 구간(믿을 만한 범위)과 함께 보여 줘요. 표~\ref{tab:tie-sensitivity}의 마지막 행은 두 점수 사이의 일치를 보여 줘요. 두 모양은 서로 엇갈려요. AWP는 송금 가족을 가장 가깝게 따라가고($0.572$, 구간 $[0.564,0.580]$), 허락 가족을 가장 덜 따라가요($0.193$, $[0.182,0.203]$). 반면에 엔도스랭크는 허락 가족을 가장 가깝게 따라가고($0.585$, $[0.578,0.592]$), 송금 가족에서는 $0.158$($[0.147,0.169]$)에 머물러요. 표~\ref{tab:tau-diff}은 계약을 다시 뽑은 같은 표본에서 짝지은 차이 다섯 개를 검사해요(\ref{sec:statistical-tests}절). 이 차이들에는 표의 행 순서대로 i부터 v까지 번호를 붙였어요.

% Generated by experiments/scripts/export_latex.py -- do not edit by hand.
% Source: experiments/data/2-processed-tables-and-evaluations/revised-usd/same-window/eval_summary.json
\begin{table}[htbp]
\centering
\caption[같은 기간에서 엔도스랭크와 AWP 사이의 짝지은 대비.]{같은 기간에서 엔도스랭크와 AWP의 켄달의 $\tau$ 계수 사이 차이를 짝지은 부트스트랩으로 검사한 결과예요(매칭 코호트, 계약 $n=15{,}107$개). 지표 가족 항목에는 지표 가족 평균 $\tau$를 써요. $a$와 $b$에 같은 계약 다시 뽑기를 써요. 그래서 구간은 짝지은 차이에 대한 구간이에요. 구간이 0을 벗어나면, 5\% 수준에서 표본을 뽑을 때 생기는 흔들림만으로는 설명할 수 없는 차이라는 뜻이에요. 신뢰구간(CI)은 계약을 다시 뽑아 구한 95\% 백분위 부트스트랩 구간이에요(짝지은 다시 뽑기 400번, 시드 42).}
\label{tab:tau-diff}
\fitwidth{%
\begin{tabular}{lccccc}
\toprule
대비 ($a$ 빼기 $b$) & $\tau_a$ & $\tau_b$ & $\Delta\tau$ & $\Delta\tau$의 95\% 신뢰구간 & $P(\Delta\tau>0)$ \\
\midrule
EndorseRank 허락 빼기 EndorseRank 송금 & 0.585 & 0.158 & +0.427 & [0.416, 0.440] & 1.000 \\
AWP 송금 빼기 EndorseRank 송금 & 0.572 & 0.158 & +0.414 & [0.401, 0.427] & 1.000 \\
EndorseRank 허락 빼기 AWP 허락 & 0.585 & 0.193 & +0.392 & [0.380, 0.405] & 1.000 \\
AWP 시빌 안정성 빼기 EndorseRank 시빌 안정성 & 0.395 & 0.007 & +0.388 & [0.376, 0.401] & 1.000 \\
EndorseRank 받은 허락 수 빼기 EndorseRank 입차수 & 0.655 & 0.182 & +0.473 & [0.461, 0.486] & 1.000 \\
\bottomrule
\end{tabular}
}
\end{table}


다섯 구간은 모두 0보다 위에 있어요. 엔도스랭크의 허락 계수는 자기의 송금 계수보다 커요(대비~i, $+0.427$, 구간 $[0.416,0.440]$). AWP의 송금 계수는 엔도스랭크의 송금 계수보다 커요(대비~ii, $+0.414$, $[0.401,0.427]$). 엔도스랭크의 허락 계수는 AWP의 허락 계수보다 커요(대비~iii, $+0.392$, $[0.380,0.405]$). AWP는 시빌 안정성 가족(시빌 공격, 곧 가짜 주소를 많이 만들어 속이는 공격을 생각해 고른 지표 묶음)에서 앞서요(대비~iv, $+0.388$, $[0.376,0.401]$). 그리고 엔도스랭크는 들어오는 화살표 수(입차수)보다 받은 허락 수를 더 가깝게 따라가요(대비~v, $+0.473$, $[0.461,0.486]$). 대비~i와~v는 대리 지표만 바꾸어 엔도스랭크를 자기 자신과 견줘요. 엔도스랭크는 계약 쌍의 $0.17\%$만 동점으로 둬요. 그래서 엔도스랭크 자신의 동점은 이 계수들 어느 것에도 천장(얻을 수 있는 가장 큰 값)을 씌우지 않아요. 그리고 대리 지표의 동점으로는 이만큼 큰 차이를 설명할 수 없어요. 입차수에서 AWP는 $0.578$($[0.569,0.587]$)에 이르지만, 엔도스랭크는 $0.182$($[0.169,0.194]$)에 이르러요. 허락 그래프 밖의 계약들을 외톨이 점(화살표가 하나도 없는 점)으로 남겨 두어도, 지표 가족 평균은 하나도 바뀌지 않아요(표~\ref{tab:tie-sensitivity}).

대비~ii와~iii의 부호가 서로 반대인 것은 구성개념에서 나와요. 각 점수는 자기 화살표로 만든 가족에서 앞서요. 그리고 이 가족들은 저마다, 그에 맞는 페이지랭크가 다듬는 입차수를 담고 있어요. 이 일치는 일부가 만든 방식 때문에 저절로 생기고, 각 점수의 의도한 구성개념을 점검해요. 이것은 점수 바깥에서 온 증거가 아니에요. 코호트의 모든 계약은 적어도 세 주인에게서 WETH, USDC, USDT0, ARB 또는 WBTC의 허락을 받았어요. 그리고 AWP의 허락 계수는 0보다 위에 있어요(표~\ref{tab:alignment-allowance}). 곧, 많은 주인이 허락한 계약은 송금 그래프에서도 중심에 있는 경향이 있어요. 어느 차이도 어느 점수가 평판을 더 잘 재는지는 말해 주지 않아요.

표~\ref{tab:robustness-sample-size}의 규모 키우기 단계들의 값은 코호트 전체의 값과 견줄 수 없어요. 이 단계들에서도 각 가족 안의 두 점수 순서는 모든 단계에서 같아요. AWP는 송금 가족과 시빌 안정성 가족에서 앞서고, 엔도스랭크는 허락 가족에서 앞서요.

\dissertationSubheading[sec:proxy-family-detail]{대리 지표 하나하나 보기}

지표 가족 평균 하나는 대리 지표마다 구한 계수 여러 개를 한데 줄인 값이에요. 표~\ref{tab:alignment-transfer}부터~\ref{tab:alignment-sybil-stability}까지는 엔도스랭크와 AWP의 계수를 대리 지표 하나씩 보고해요. 스피어먼의 로 $\rho$(두 줄 세우기가 얼마나 비슷한지 나타내는 수), 켄달의 타우 $\tau$(두 줄 세우기가 얼마나 같은 순서인지 나타내는 수), 그리고 $\tau$의 95\% 구간이에요. 모든 대리 지표는 다섯 토큰의 이벤트(계약이 남기는 기록)를 세고, 금액은 달러로 매겨요(\ref{sec:proxy-formulas}절). 상관을 구하기 전에, 모든 대리 지표는 값이 클수록 기대되는 평판이 높다는 뜻이 되도록 방향을 맞춰요. 그리고 송금 대리 지표는 AWP를 만든 사람들이 AWP를 평가할 때 쓴 바로 그 지표예요\parencite{do2023awp}.

\dissertationSubsubheading[sub:transfer-family]{송금 가족}

표~\ref{tab:alignment-transfer}은 두 점수와 입차수, 그리고 받은 양의 합 사이의 계수를 보여 줘요.

% Generated by experiments/scripts/export_latex.py -- do not edit by hand.
% Source: experiments/data/2-processed-tables-and-evaluations/revised-usd/same-window/eval_summary.json
\begin{table}[htbp]
\centering
\caption[송금 대리 지표에 대한 엔도스랭크와 AWP의 분야 일치도.]{엔도스랭크와 AWP가 송금 대리 지표(각 계약의 입차수와 받은 양의 합)와 얼마나 잘 맞는지 나타낸 분야 일치도예요. 같은 기간에서 쟀어요(매칭 코호트, 계약 $n=15{,}107$개). 점수의 순위와 대리 지표의 순위 사이의 켄달의 $\tau$와 스피어먼의 $\rho$를 보여 줘요. 신뢰구간(CI)은 계약을 다시 뽑아 구한 95\% 백분위 부트스트랩 구간이에요(짝지은 다시 뽑기 400번, 시드 42).}
\label{tab:alignment-transfer}
\fitwidth{%
\begin{tabular}{lccc}
\toprule
대리 지표 & 스피어먼의 $\rho$ & 켄달의 $\tau$ & $\tau$의 95\% 신뢰구간 \\
\midrule
\multicolumn{4}{l}{\textit{EndorseRank}} \\
입차수 & 0.261 & 0.182 & [0.169, 0.194] \\
받은 양의 합 (들어온 양을 모두 더함) & 0.205 & 0.134 & [0.123, 0.145] \\
\addlinespace
\multicolumn{4}{l}{\textit{AWP}} \\
입차수 & 0.744 & 0.578 & [0.569, 0.587] \\
받은 양의 합 (들어온 양을 모두 더함) & 0.741 & 0.566 & [0.558, 0.575] \\
\bottomrule
\end{tabular}
}
\end{table}


AWP는 입차수와 $0.578$(구간 $[0.569,0.587]$)만큼 맞고, 받은 양의 합과 $0.566$($[0.558,0.575]$)만큼 맞아요. AWP의 금액 무게는 중앙값 송금인 200달러에서 2분의 1이에요. 그리고 그보다 몇 배 큰 금액에서만 1에 가까워져요(\ref{sub:awp-graph}절). 그래서 큰 송금일수록 더 크게 세요. 그리고 AWP는 한 계약이 받은 송금의 달러 합계를, 그 송금의 수와 거의 같은 만큼 가깝게 따라가요. 엔도스랭크의 송금 계수는 $0.182$($[0.169,0.194]$)와 $0.134$($[0.123,0.145]$)로 0보다 위에 있어요. 하지만 그 구간은 AWP의 구간보다 훨씬 아래에 있어요. 곧, 많은 주인이 허락한 계약은 코호트의 나머지 계약보다 더 많은 보낸 쪽(송신자)에게서 송금을 받아요. 기간 밖에서 AWP는 두 기간 모두에서 나중의 새 송금자를 원시 입차수보다 더 잘 줄 세워요(\ref{sec:holdout-results}절). 그러니 송금 쪽에서는 걷기가 자기가 다듬는 수에 무언가를 보태요.

\dissertationSubsubheading[sub:allowance-family]{허락 가족}

표~\ref{tab:alignment-allowance}는 두 점수와 받은 허락 수, 그리고 받은 허락 금액 사이의 계수를 보여 줘요.

% Generated by experiments/scripts/export_latex.py -- do not edit by hand.
% Source: experiments/data/2-processed-tables-and-evaluations/revised-usd/same-window/eval_summary.json
\begin{table}[htbp]
\centering
\caption[허락 대리 지표에 대한 엔도스랭크와 AWP의 분야 일치도.]{엔도스랭크와 AWP가 허락 대리 지표(각 스펜더 계약의 받은 허락 수와 받은 허락의 양)와 얼마나 잘 맞는지 나타낸 분야 일치도예요. 같은 기간에서 쟀어요(매칭 코호트, 계약 $n=15{,}107$개). 점수의 순위와 대리 지표의 순위 사이의 켄달의 $\tau$와 스피어먼의 $\rho$를 보여 줘요. 신뢰구간(CI)은 계약을 다시 뽑아 구한 95\% 백분위 부트스트랩 구간이에요(짝지은 다시 뽑기 400번, 시드 42).}
\label{tab:alignment-allowance}
\fitwidth{%
\begin{tabular}{lccc}
\toprule
대리 지표 & 스피어먼의 $\rho$ & 켄달의 $\tau$ & $\tau$의 95\% 신뢰구간 \\
\midrule
\multicolumn{4}{l}{\textit{EndorseRank}} \\
받은 허락 수 & 0.811 & 0.655 & [0.648, 0.662] \\
받은 허락의 양 (합) & 0.681 & 0.515 & [0.506, 0.524] \\
\addlinespace
\multicolumn{4}{l}{\textit{AWP}} \\
받은 허락 수 & 0.286 & 0.206 & [0.193, 0.217] \\
받은 허락의 양 (합) & 0.260 & 0.180 & [0.169, 0.191] \\
\bottomrule
\end{tabular}
}
\end{table}


엔도스랭크는 받은 허락 수와 $0.655$(구간 $[0.648,0.662]$)만큼 맞아요. 그리고 받은 허락 금액과는 $0.515$($[0.506,0.524]$)만큼 맞아요. 받은 허락 금액은 상한(이보다 큰 금액은 이 값으로 자름)으로 자른 허락들을 더한 값이에요. 금액 무게는 끝이 있는 금액의 중앙값인 62달러짜리 허락에 2분의 1을 줘요. 하지만 0보다 큰 마지막 허락(가장 최근 허락) 가운데 절반이 넘는 것이 무제한 허락이거나 20{,}000달러에 이르러요(표~\ref{tab:usd-tokens}). 이런 허락은 꽉 채워 세요. 그래서 엔도스랭크는 주인들이 허락한 금액의 합보다, 허락한 주인의 수를 더 가깝게 따라가요. 주인들이 허락을 하나도 받지 않은 스펜더에서는, 엔도스랭크가 \ref{sub:theory-closed-form}절의 닫힌 식(바로 계산할 수 있는 식)과 같아요. 이 식은 주인 정규화 입강도(각 주인이 준 무게를 그 주인이 허락한 모든 스펜더에게 나누어 더한 값)에 따라 정해져요. AWP는 두 대리 지표와 $0.206$($[0.193,0.217]$), $0.180$($[0.169,0.191]$)만큼 맞아요. 이 비교는 일부가 만든 방식 때문에 저절로 그렇게 돼요. 엔도스랭크가 대리 지표와 같은 허락들로 계산되기 때문이에요.

\dissertationSubsubheading[sub:sybil-family]{시빌 보정 안정성 가족}

표~\ref{tab:alignment-sybil-stability}은 두 점수와 들어오는 상대 비율, 송금 활동 기간, 활동한 달 수 사이의 계수를 보여 줘요.

% Generated by experiments/scripts/export_latex.py -- do not edit by hand.
% Source: experiments/data/2-processed-tables-and-evaluations/revised-usd/same-window/eval_summary.json
\begin{table}[htbp]
\centering
\caption[시빌 보정 안정성 대리 지표에 대한 엔도스랭크와 AWP의 분야 일치도.]{엔도스랭크와 AWP가 송금 활동에서 구한 시빌 보정 안정성 대리 지표와 얼마나 잘 맞는지 나타낸 분야 일치도예요. 같은 기간에서 쟀어요(매칭 코호트, 계약 $n=15{,}107$개). 점수의 순위와 대리 지표의 순위 사이의 켄달의 $\tau$와 스피어먼의 $\rho$를 보여 줘요. 신뢰구간(CI)은 계약을 다시 뽑아 구한 95\% 백분위 부트스트랩 구간이에요(짝지은 다시 뽑기 400번, 시드 42).}
\label{tab:alignment-sybil-stability}
\fitwidth{%
\begin{tabular}{lccc}
\toprule
대리 지표 & 스피어먼의 $\rho$ & 켄달의 $\tau$ & $\tau$의 95\% 신뢰구간 \\
\midrule
\multicolumn{4}{l}{\textit{EndorseRank}} \\
들어오는 상대 비율 & -0.231 & -0.156 & [-0.166, -0.147] \\
송금 활동 기간 (일) & 0.124 & 0.080 & [0.069, 0.091] \\
활동한 달 수 & 0.146 & 0.098 & [0.087, 0.109] \\
\addlinespace
\multicolumn{4}{l}{\textit{AWP}} \\
들어오는 상대 비율 & 0.234 & 0.127 & [0.114, 0.141] \\
송금 활동 기간 (일) & 0.683 & 0.508 & [0.499, 0.517] \\
활동한 달 수 & 0.704 & 0.549 & [0.540, 0.558] \\
\bottomrule
\end{tabular}
}
\end{table}


표를 보면 AWP는 송금 활동 기간에서 엔도스랭크보다 앞서요($0.508$, 구간 $[0.499,0.517]$, 대 $0.080$, $[0.069,0.091]$). 활동한 달 수에서도 앞서요($0.549$, $[0.540,0.558]$, 대 $0.098$, $[0.087,0.109]$). 기간 전체에 퍼진 활동에 점수를 더 주는 송금 점수라면 그럴 것이라고 예상할 수 있어요. 지표 가족 평균에서의 짝지은 차이는 0을 벗어나요(표~\ref{tab:tau-diff}, 대비~iv). 들어오는 상대 비율에서는 두 계수의 부호가 서로 반대예요. 엔도스랭크는 $-0.156$($[-0.166,-0.147]$)이고, AWP는 $0.127$($[0.114,0.141]$)이에요. 이 비율은 자기 자신을 뺀 서로 다른 보낸 쪽의 수를, 그들에게서 받은 송금의 수로 나눈 값이에요. 그래서 같은 보낸 쪽이 다시 오면 비율이 내려가요. 엔도스랭크의 계수가 음수라는 것은, 많은 주인이 허락한 계약이 이 잣대로 보면 보낸 쪽이 다시 찾아오는 계약이라는 뜻이에요. AWP의 계수가 양수인 것은 점수를 매기는 범위 때문이에요. 송금 그래프 밖에 있는 계약 $2{,}670$개는 비율도 0이고 AWP 점수도 0이에요. 그래서 이 계약들이 두 줄 세우기의 맨 아래를 함께 채워요. AWP가 0보다 큰 점수를 준 계약들만 보면, AWP의 순서도 엔도스랭크처럼 비율이 커질수록 내려가요. 이 점검은 결과를 안 뒤에 했고, 구간은 없어요. 이 가족은 각 점수가 시간에 걸친 송금 정보를 얼마나 담고 있는지를 보여 줘요.

\dissertationSubheading[sec:er-awp-holdout]{기간 밖}

표~\ref{tab:endorserank-awp-holdout}은 두 홀드아웃(나중 기간을 떼어 두고 맞혀 보는 시험) 기간에서 두 점수 사이의 짝지은 차이를 모아요. 모든 행에서 $a$는 엔도스랭크이고 $b$는 AWP예요.

% Generated by experiments/scripts/export_latex.py -- do not edit by hand.
% Source: experiments/data/2-processed-tables-and-evaluations/revised-usd/holdout-w0/eval_summary.json
% Source: experiments/data/2-processed-tables-and-evaluations/revised-usd/supplementary/supplementary.json
% Source: experiments/data/2-processed-tables-and-evaluations/revised-usd/window-w1/eval_summary.json
\begin{table}[htbp]
\centering
\caption[기간 밖에서 엔도스랭크와 AWP 사이의 짝지은 부트스트랩 차이.]{기간 밖에서 엔도스랭크와 AWP 사이의 차이를 짝지은 부트스트랩으로 구했어요(W0 스펜더 $n=17{,}050$, W0 트레이더 $n=88$, W1 스펜더 $n=18{,}152$, W1 트레이더 $n=70$). 이 가운데 어느 것도 규칙 A나 규칙 B에 들어가지 않아요. 스펜더 행은 라벨을 본 뒤에 계산했어요. 트레이더 행은 어느 화살표 종류로도 만들지 않은 결과에서 한 비교 가운데, 엔도스랭크에서 AWP를 뺀 쌍이에요(표~\ref{tab:neutral-label-contrasts}과~\ref{tab:neutral-label-grid}). W1 스펜더 행은 W1 평가가 끝난 뒤에, 저장해 둔 W1 점수로 계산했어요. 다시 뽑은 표본은 표~\ref{tab:fresh-holdout}와 같은 것을 썼어요. 각 묶음 안에서 $a$와 $b$는 같은 계약 다시 뽑은 표본을 함께 써요(스펜더는 400번, 트레이더는 2{,}000번, 시드 42).}
\label{tab:endorserank-awp-holdout}
\fitwidth{%
\begin{tabular}{llccccc}
\toprule
대비 ($a$ 빼기 $b$) & 라벨 & $\tau_a$ & $\tau_b$ & $\Delta\tau$ & $\Delta\tau$의 95\% 신뢰구간 & $P(\Delta\tau>0)$ \\
\midrule
\multicolumn{7}{l}{\emph{W0, 스펜더 ($n=17{,}050$)}} \\
EndorseRank 빼기 AWP & 새 허락 쌍 & 0.231 & 0.196 & +0.035 & [0.019, 0.051] & 1.000 \\
EndorseRank 빼기 AWP & 새 송금자 & 0.200 & 0.346 & -0.146 & [-0.156, -0.136] & 0.000 \\
\midrule
\multicolumn{7}{l}{\emph{W0, 세 번 이상 닫은 트레이더 ($n=88$)}} \\
EndorseRank 빼기 AWP & 청산 없이 닫은 비율 & 0.224 & 0.240 & -0.017 & [-0.156, 0.128] & 0.394 \\
\midrule
\multicolumn{7}{l}{\emph{W1, 스펜더 ($n=18{,}152$)}} \\
EndorseRank 빼기 AWP & 새 허락 쌍 & 0.234 & 0.169 & +0.065 & [0.051, 0.081] & 1.000 \\
EndorseRank 빼기 AWP & 새 송금자 & 0.194 & 0.330 & -0.136 & [-0.146, -0.127] & 0.000 \\
\midrule
\multicolumn{7}{l}{\emph{W1, 세 번 이상 닫은 트레이더 ($n=70$)}} \\
EndorseRank 빼기 AWP & 청산 없이 닫은 비율 & 0.297 & 0.329 & -0.032 & [-0.133, 0.066] & 0.279 \\
\bottomrule
\end{tabular}
}
\end{table}


W0의 스펜더에서 엔도스랭크는 새 허락 쌍을 $\tau=0.231$(구간 $[0.220,0.242]$)로 줄 세워요. AWP는 $0.196$($[0.181,0.210]$)이에요. 짝지은 차이는 $+0.035$($[0.019,0.051]$)예요. 새 송금자에서는 순서가 뒤집혀요. AWP는 $0.346$($[0.336,0.356]$)에 이르고, 엔도스랭크는 $0.200$($[0.189,0.210]$)이에요. 차이는 $-0.146$($[-0.156,-0.136]$)이에요. 모든 라벨은 점수에 쓴 모든 이벤트보다 뒤에 와요. 그러니 이 코호트에서는 각 점수가 자기 화살표 종류로 나중에 들어온 것을 다른 점수보다 더 잘 줄 세워요. 첫 번째 차이의 일부는 점수를 매기는 범위 때문이에요. AWP는 동결 전에 송금 화살표가 없던 스펜더 $2{,}891$개에 0점을 줘요. 그런데 이 스펜더의 $13.3\%$가 새 허락 쌍을 얻었어요. 나머지 스펜더는 $11.9\%$였어요. 다른 주소에게서 송금을 받은 적이 있는 스펜더 $14{,}094$개에서는, 새 허락 쌍에서의 순서가 뒤집혀요. AWP는 $0.298$($[0.287,0.308]$)에 이르고, 엔도스랭크는 $0.236$($[0.224,0.249]$)이에요. 여기에는 짝지은 검사를 하지 않았어요(표~\ref{tab:holdout-receiving}). 이 스펜더들에서도 새 송금자로는 AWP가 계속 앞서요. $0.357$($[0.347,0.367]$) 대 $0.231$($[0.219,0.244]$)이에요.

7월부터 9월까지의 기간에서도, 스펜더 코호트 전체에서 같은 모양이 되풀이돼요. 이때 대비는 모든 라벨을 본 뒤에, 수정한 W1 점수로 계산했어요. 엔도스랭크는 새 허락 쌍을 $0.234$($[0.222,0.244]$)로 줄 세우고, AWP는 $0.169$($[0.157,0.183]$)예요($+0.065$, $[0.051,0.081]$). AWP는 새 송금자를 $0.330$($[0.321,0.340]$)으로 줄 세우고, 엔도스랭크는 $0.194$($[0.184,0.205]$)예요($-0.136$, $[-0.146,-0.127]$). 트레이더의 청산 없이 닫은 비율(강제로 끝나지 않고 거래를 닫은 비율)도 보아요. 이 쌍은 W1의 등록 문서가 등록 분석을 위해 정해 둔 쌍이에요. 엔도스랭크는 W0에서 $0.224$($[0.105,0.339]$), W1에서 $0.297$($[0.171,0.410]$)에 이르러요. AWP는 $0.240$($[0.073,0.380]$)과 $0.329$($[0.207,0.440]$)예요. 두 차이 모두 0과 가려낼 수 없어요($-0.017$, $[-0.156,0.128]$, 그리고 $-0.032$, $[-0.133,0.066]$). 트레이더 코호트에는 계약이 각각 $88$개와 $70$개 있어서, 구간이 넓어요.

W0 트레이더의 이익 라벨(표~\ref{tab:holdout-traders})에서, 이익을 낸 닫기와 실현 이익(실제로 얻은 이익)에 대한 엔도스랭크의 구간은 0을 품어요. 값은 $0.150$($[-0.019,0.295]$)과 $0.078$($[-0.094,0.226]$)이에요. 반면에 AWP의 구간은 0보다 위에 있어요. 값은 $0.155$($[0.005,0.297]$)와 $0.196$($[0.048,0.351]$)이에요. 두 라벨은 모두 닫기 횟수가 많을수록 커져요. 두 비율은 그렇지 않은데, 여기서는 두 점수가 모두 음수예요. 그리고 두 점수 모두 손실이 아닌 비율에서만 구간이 0보다 아래에 있어요. AWP는 손실이 아닌 비율에서 $-0.254$($[-0.404,-0.104]$)이고, 이익을 낸 비율에서는 $-0.124$($[-0.265,0.015]$)예요. 엔도스랭크는 $-0.178$($[-0.298,-0.042]$) 대 $-0.055$($[-0.196,0.084]$)예요. 송금 입차수도 두 비율 모두에서 음수예요. 값은 $-0.191$($[-0.329,-0.041]$)과 $-0.197$($[-0.350,-0.035]$)이에요. 어느 점수도 이 계약 트레이더(GMX V2에서 거래하는 계약)들을 성공한 닫기의 비율대로 줄 세우지 못해요.

\dissertationSubheading[sec:er-awp-neutral-grid]{두 기간의 모든 트레이더 라벨}

표~\ref{tab:neutral-label-grid}는 두 기간에서 트레이더 라벨 여섯 개 모두에 대해 허락 쪽과 송금 쪽을 맞세워요. 이 표의 쌍 네 개와 라벨 여섯 개는 등록 분석을 위해 W1과 함께 등록했어요. 7월부터 9월까지의 로그를 뽑기 전이었어요(\ref{sub:neutral-label-protocol}절). 그때까지 등록 분석은 W0 트레이더를 이익 라벨로는 채점했지만, 청산 라벨로는 채점하지 않았어요. 이 표는 같은 쌍을 고친 점수로 보여 줘요. 고친 점수는 청산 라벨을 포함한 모든 등록 분석의 결과를 안 뒤에 정의했어요. 그래서 이 표의 값들은 재분석이지, 등록한 시험이 아니에요(부록~\ref{sec:map-registered}). 청산 없이 닫은 비율에서의 P와 Q는, W1의 값과 W0의 부호로 규칙 R1과 R2에 들어가요. W1의 두 비율에서의 P는 규칙 R3에 들어가요. 이 규칙들은 모두 고친 점수에 적용해요(주장~C6, \ref{sub:neutral-label-results}절). 나머지 값들은 엔도스랭크 빼기 AWP 열까지 포함해서, 그 옆에 함께 보고해요.

% Generated by experiments/scripts/export_latex.py -- do not edit by hand.
% Source: experiments/data/2-processed-tables-and-evaluations/revised-usd/window-w1/eval_summary.json
\begin{table}[htbp]
\centering
\caption[두 기간의 모든 트레이더 라벨에서 허락 쪽 점수 빼기 송금 쪽 점수.]{허락 쪽 점수에서 송금 쪽 점수를 뺀 $\Delta\tau$와 95\% 짝지은 부트스트랩 구간이에요. 기간 W1(GMX~V2 계약 트레이더 $n=70$)과 W0($n=88$)의 트레이더 라벨 여섯 개 모두에서 구했어요. P는 받은 허락 수 빼기 송금 입차수예요. Q는 활동 재시작을 쓰는 엔도스랭크 빼기 AWP예요(화살표 무게와 다시 출발 규칙이 같음). C-PR은 $\lambda=1$ 빼기 $\lambda=0$이에요(같은 달러 방식으로 만듦). ER은 엔도스랭크 빼기 AWP예요. 이익을 낸 닫기 수와 실현 이익은 닫은 횟수가 많을수록 커져요. 다른 네 라벨은 그렇지 않아요. 신뢰구간(CI)은 계약을 다시 뽑아 구한 95\% 백분위 부트스트랩 구간이에요(짝지은 다시 뽑기 2{,}000번, 시드 42).}
\label{tab:neutral-label-grid}
\fitwidth{%
\begin{tabular}{lcccc}
\toprule
라벨 & P (개수) & Q (페이지랭크) & C-PR 층 & ER 빼기 AWP \\
\midrule
\multicolumn{5}{l}{\emph{기간 W1, 2026년 7월--9월 (트레이더 $n=70$)}} \\
청산 없이 닫은 비율 & \begin{tabular}[c]{@{}c@{}}0.084\\[-1pt]{\footnotesize [-0.013, 0.176]}\end{tabular} & \begin{tabular}[c]{@{}c@{}}-0.041\\[-1pt]{\footnotesize [-0.135, 0.050]}\end{tabular} & \begin{tabular}[c]{@{}c@{}}-0.049\\[-1pt]{\footnotesize [-0.169, 0.060]}\end{tabular} & \begin{tabular}[c]{@{}c@{}}-0.032\\[-1pt]{\footnotesize [-0.133, 0.066]}\end{tabular} \\
청산 없음 & \begin{tabular}[c]{@{}c@{}}0.078\\[-1pt]{\footnotesize [-0.012, 0.166]}\end{tabular} & \begin{tabular}[c]{@{}c@{}}-0.039\\[-1pt]{\footnotesize [-0.138, 0.062]}\end{tabular} & \begin{tabular}[c]{@{}c@{}}-0.051\\[-1pt]{\footnotesize [-0.170, 0.064]}\end{tabular} & \begin{tabular}[c]{@{}c@{}}-0.029\\[-1pt]{\footnotesize [-0.136, 0.080]}\end{tabular} \\
이익을 낸 닫기 수 & \begin{tabular}[c]{@{}c@{}}-0.143\\[-1pt]{\footnotesize [-0.318, 0.013]}\end{tabular} & \begin{tabular}[c]{@{}c@{}}-0.231\\[-1pt]{\footnotesize [-0.429, -0.027]}\end{tabular} & \begin{tabular}[c]{@{}c@{}}-0.190\\[-1pt]{\footnotesize [-0.411, 0.037]}\end{tabular} & \begin{tabular}[c]{@{}c@{}}-0.214\\[-1pt]{\footnotesize [-0.423, -0.012]}\end{tabular} \\
실현 이익 & \begin{tabular}[c]{@{}c@{}}-0.135\\[-1pt]{\footnotesize [-0.308, 0.025]}\end{tabular} & \begin{tabular}[c]{@{}c@{}}-0.350\\[-1pt]{\footnotesize [-0.536, -0.157]}\end{tabular} & \begin{tabular}[c]{@{}c@{}}-0.404\\[-1pt]{\footnotesize [-0.624, -0.182]}\end{tabular} & \begin{tabular}[c]{@{}c@{}}-0.346\\[-1pt]{\footnotesize [-0.547, -0.144]}\end{tabular} \\
이익을 낸 비율 & \begin{tabular}[c]{@{}c@{}}-0.196\\[-1pt]{\footnotesize [-0.364, -0.046]}\end{tabular} & \begin{tabular}[c]{@{}c@{}}-0.336\\[-1pt]{\footnotesize [-0.523, -0.141]}\end{tabular} & \begin{tabular}[c]{@{}c@{}}-0.344\\[-1pt]{\footnotesize [-0.551, -0.120]}\end{tabular} & \begin{tabular}[c]{@{}c@{}}-0.316\\[-1pt]{\footnotesize [-0.517, -0.112]}\end{tabular} \\
손실이 아닌 비율 & \begin{tabular}[c]{@{}c@{}}-0.252\\[-1pt]{\footnotesize [-0.391, -0.120]}\end{tabular} & \begin{tabular}[c]{@{}c@{}}-0.350\\[-1pt]{\footnotesize [-0.526, -0.175]}\end{tabular} & \begin{tabular}[c]{@{}c@{}}-0.324\\[-1pt]{\footnotesize [-0.531, -0.109]}\end{tabular} & \begin{tabular}[c]{@{}c@{}}-0.346\\[-1pt]{\footnotesize [-0.526, -0.161]}\end{tabular} \\
\midrule
\multicolumn{5}{l}{\emph{기간 W0, 2026년 4월--6월 (트레이더 $n=88$)}} \\
청산 없이 닫은 비율 & \begin{tabular}[c]{@{}c@{}}0.182\\[-1pt]{\footnotesize [0.052, 0.314]}\end{tabular} & \begin{tabular}[c]{@{}c@{}}-0.014\\[-1pt]{\footnotesize [-0.147, 0.130]}\end{tabular} & \begin{tabular}[c]{@{}c@{}}-0.041\\[-1pt]{\footnotesize [-0.187, 0.105]}\end{tabular} & \begin{tabular}[c]{@{}c@{}}-0.017\\[-1pt]{\footnotesize [-0.156, 0.128]}\end{tabular} \\
청산 없음 & \begin{tabular}[c]{@{}c@{}}0.197\\[-1pt]{\footnotesize [0.059, 0.333]}\end{tabular} & \begin{tabular}[c]{@{}c@{}}-0.010\\[-1pt]{\footnotesize [-0.152, 0.146]}\end{tabular} & \begin{tabular}[c]{@{}c@{}}-0.051\\[-1pt]{\footnotesize [-0.198, 0.102]}\end{tabular} & \begin{tabular}[c]{@{}c@{}}-0.012\\[-1pt]{\footnotesize [-0.154, 0.137]}\end{tabular} \\
이익을 낸 닫기 수 & \begin{tabular}[c]{@{}c@{}}0.075\\[-1pt]{\footnotesize [-0.080, 0.222]}\end{tabular} & \begin{tabular}[c]{@{}c@{}}-0.016\\[-1pt]{\footnotesize [-0.150, 0.118]}\end{tabular} & \begin{tabular}[c]{@{}c@{}}-0.067\\[-1pt]{\footnotesize [-0.222, 0.087]}\end{tabular} & \begin{tabular}[c]{@{}c@{}}-0.005\\[-1pt]{\footnotesize [-0.163, 0.155]}\end{tabular} \\
실현 이익 & \begin{tabular}[c]{@{}c@{}}0.017\\[-1pt]{\footnotesize [-0.150, 0.176]}\end{tabular} & \begin{tabular}[c]{@{}c@{}}-0.107\\[-1pt]{\footnotesize [-0.252, 0.044]}\end{tabular} & \begin{tabular}[c]{@{}c@{}}-0.248\\[-1pt]{\footnotesize [-0.433, -0.067]}\end{tabular} & \begin{tabular}[c]{@{}c@{}}-0.118\\[-1pt]{\footnotesize [-0.286, 0.046]}\end{tabular} \\
이익을 낸 비율 & \begin{tabular}[c]{@{}c@{}}0.045\\[-1pt]{\footnotesize [-0.113, 0.184]}\end{tabular} & \begin{tabular}[c]{@{}c@{}}0.046\\[-1pt]{\footnotesize [-0.071, 0.172]}\end{tabular} & \begin{tabular}[c]{@{}c@{}}0.027\\[-1pt]{\footnotesize [-0.112, 0.181]}\end{tabular} & \begin{tabular}[c]{@{}c@{}}0.069\\[-1pt]{\footnotesize [-0.060, 0.208]}\end{tabular} \\
손실이 아닌 비율 & \begin{tabular}[c]{@{}c@{}}-0.080\\[-1pt]{\footnotesize [-0.228, 0.058]}\end{tabular} & \begin{tabular}[c]{@{}c@{}}0.055\\[-1pt]{\footnotesize [-0.070, 0.192]}\end{tabular} & \begin{tabular}[c]{@{}c@{}}0.037\\[-1pt]{\footnotesize [-0.111, 0.189]}\end{tabular} & \begin{tabular}[c]{@{}c@{}}0.076\\[-1pt]{\footnotesize [-0.063, 0.223]}\end{tabular} \\
\bottomrule
\end{tabular}
}
\end{table}


W1에서는 어느 쌍도 두 청산 라벨 어디에서든 두 쪽을 갈라놓지 못해요. 모든 구간이 0을 품어요. 청산 없이 닫은 비율에서 P는 $+0.084$($[-0.013,0.176]$)이고, 엔도스랭크 쌍은 $-0.032$($[-0.133,0.066]$)예요. W0에서는 두 청산 라벨 모두에서 그냥 센 수가 허락 쪽에 유리해요. P는 청산 없이 닫은 비율에서 $+0.182$($[0.052,0.314]$), 청산이 없었는지 나타내는 표시에서 $+0.197$($[0.059,0.333]$)이에요. 반면에 페이지랭크 쌍 세 개는 0보다 아래에 있고, 그 구간은 0을 품어요. W1의 이익 라벨 네 개에서는 모든 쌍이 음수예요. Q와 엔도스랭크 쌍의 구간은 네 라벨 모두에서 0보다 아래에 있어요. P의 구간은 두 비율에서, C-PR(결합 페이지랭크, 두 그래프를 함께 걷는 점수) 층들의 구간은 실현 이익과 두 비율에서 0보다 아래에 있어요. W0에서는 이익 라벨에서의 모든 구간이 0을 품어요. 하나만 예외예요. 실현 이익에서의 C-PR 층들의 구간이고, 값은 $-0.248$($[-0.433,-0.067]$)이에요. 그래서 실현 이익에서의 C-PR 층들의 구간은 두 기간 모두에서 0보다 아래에 있어요. 실현 이익은 닫기 횟수가 많을수록 커지는 라벨이에요. W0의 두 비율에서는 Q, C-PR 층들, 엔도스랭크 쌍이 허락 쪽으로 기울지만, 그 구간은 0을 품어요. 그래서 두 비율은 두 기간 모두에서 유지되는 방향을 보여 주지 않아요. 트레이더가 $70$개와 $88$개라서, 이 차이들 가운데 어느 것도 정밀하지 않아요.
